
Note also that, since @@M0@@ is regular and hence normal, its connected components coincide with its irreducible components, which dominate @@M1@@, so that the homomorphism @@M2@@ induced by the inclusion @@M3@@ is bijective. Consequently the homomorphism @@M4@@ induced by the inclusion @@M5@@ can also be interpreted as a homomorphism
\sgaStableTarget{sga6.e10.eq.7-7-3}@@M6@@
which again deserves the name ``specialization homomorphism''. One then checks immediately from the definitions that the homomorphism (7.7.1) is compatible with the \emph{augmented algebra} structures over @@M7@@ and @@M8@@ respectively on the two sides, defined in (IV 2.8), taking account of the preceding homomorphism (7.7.3).

Still supposing that @@M9@@ is separated, the preceding discussion implies in particular that (7.7.1) induces a homomorphism of graded rings
\sgaStableTarget{sga6.e10.eq.7-7-4}@@M10@@
and, using the notion of Chern class of (V 6), one obtains a commutative diagram
@@M11@@

\sgaStableTarget{sga6.e10.sec.7.8}\subsection*{7.8. Specialization for @@M12@@.}
Retain the hypotheses of 7.7: @@M13@@ separated, @@M14@@ and @@M15@@ regular. Taking account of the canonical isomorphism (5.3.2):
@@M16@@
the homomorphism (7.7.4) therefore gives a canonical homomorphism
\sgaStableTarget{sga6.e10.eq.7-8-1}@@M17@@
